% !TEX root = ../main.tex
\section{引言}
\label{sec:intro}

\planned{按由现象逐步收束到机制、再展开到方法与证据的 V 形结构组织，以下各写一段。
(1) 现象：推理智能体完成一项任务需要数分钟；几乎所有时间都花在目标模型自身的生成上，且各轮次串行执行。
(2) 现有技术的局限：工具重叠推测执行~\citep{bai2026spork,nichols2025speculative} 的收益受工具耗时占比限制；跨步骤推测执行~\citep{ye2026speculative,xu2026leap} 必须在目标模型完成推理前预测其决策，因此需要针对工作负载训练的草拟器，或接受较低的一致率。
(3) 机制：目标模型在自身推理过程中就会显露决策；当历史不保留推理内容时，下一轮输入对当前轮的依赖仅来自可见输出和观测。循环等待的是确认，而非新信息。
(4) 方法：\method 及其三项策略。
(5) 证据与贡献：在实验部分完成后确定。}
